% ==============================================================================
% PREAMBLE & KONFIGURASI PACKAGES
% ==============================================================================

% Pengaturan Halaman & Margin (A4, standar 4-3-3-3 atau 3-3-3-3 cm)
\usepackage[a4paper, left=3cm, right=3cm, top=3cm, bottom=3cm]{geometry}

% Dukungan Bahasa & Karakter
\usepackage[utf8]{inputenc}
\usepackage[indonesian]{babel} % Menggunakan Bahasa Indonesia untuk penamaan (Gambar, Tabel, Bab)

% Font & Tipografi
\usepackage{times} % Menggunakan font Times New Roman (standar akademik)
\usepackage{microtype}
\usepackage{setspace}
\onehalfspacing % Spasi 1.5 (standar makalah)

% Grafik & Gambar
\usepackage{graphicx}
\usepackage{float} % Memaksa posisi gambar dengan [H]
\usepackage{caption}
\usepackage{subcaption}
\graphicspath{{figures/}} % Lokasi folder gambar

% Tabel
\usepackage{booktabs} % Garis tabel standar profesional (toprule, midrule, bottomrule)
\usepackage{multirow}
\usepackage{array}

% Matematika
\usepackage{amsmath, amssymb, amsfonts, amsthm}

% Source Code (Kodingan Robot/Data)
\usepackage{listings}
\usepackage{xcolor}

% Konfigurasi Tampilan Source Code
\definecolor{codegreen}{rgb}{0,0.6,0}
\definecolor{codegray}{rgb}{0.5,0.5,0.5}
\definecolor{codepurple}{rgb}{0.58,0,0.82}
\definecolor{backcolour}{rgb}{0.95,0.95,0.92}

\lstdefinestyle{robotstyle}{
    backgroundcolor=\color{backcolour},   
    commentstyle=\color{codegreen},
    keywordstyle=\color{blue},
    numberstyle=\tiny\color{codegray},
    stringstyle=\color{codepurple},
    basicstyle=\ttfamily\footnotesize,
    breakatwhitespace=false,         
    breaklines=true,                 
    captionpos=b,                    
    keepspaces=true,                 
    numbers=left,                    
    numbersep=5pt,                  
    showspaces=false,                
    showstringspaces=false,
    showtabs=false,                  
    tabsize=2
}
\lstset{style=robotstyle}

% Dukungan Kutipan & Bibliografi (Numeric Style standar)
\usepackage{csquotes} % Direkomendasikan untuk babel/polyglossia
\DeclareQuoteAlias{english}{indonesian}
\usepackage[backend=biber, style=numeric, sorting=none]{biblatex}
\DeclareLanguageMapping{indonesian}{english}
\addbibresource{references.bib}

% Hyperlinks (Link aktif untuk daftar isi, gambar, referensi)
\usepackage{hyperref}
\hypersetup{
    hypertexnames=false,
    colorlinks=true,
    linkcolor=black,
    filecolor=magenta,      
    urlcolor=blue,
    pdftitle={Makalah UAS ISR},
    pdfpagemode=FullScreen,
}

% Pengaturan Paragraf & Indentasi
\setlength{\parindent}{1cm}
\setlength{\parskip}{0.2cm}
